\section{The reference map and the elastic solid}
\label{sec:rmt}

\subsection{Reference map and deformation gradient}
Let $\bm\chi(\Xv,t)$ be the motion carrying a material label $\Xv$ to its current
position $\xv$. The \emph{reference map} is the inverse,
$\xiv(\xv,t)=\bm\chi^{-1}(\xv,t)$, so $\xiv(\xv,0)=\xv$ in the solid. Because $\xiv$
labels material, it is advected passively,
\begin{equation}
\partial_t\xiv + (\uv\!\cdot\!\grad)\xiv = \bm 0 .
\label{eq:xi-advect}
\end{equation}
The deformation gradient $\Ften=\partial\xv/\partial\Xv$ is recovered in the Eulerian
frame from the spatial gradient of the reference map,
\begin{equation}
\Ften = (\grad\xiv)^{-1},
\qquad J=\det\Ften ,
\label{eq:F}
\end{equation}

\phfig{fig:rmt-schematic}{Reference map technique. \TODO{schematic: (left) material body
with label field $\Xv$ in the reference configuration; (centre) the motion $\bm\chi$ and
its inverse, the reference map $\xiv=\bm\chi^{-1}$, on the fixed Eulerian grid; (right)
$\Ften=(\grad\xiv)^{-1}$ and the solid/fluid blend across the smoothed interface $\phi$.
Draw as inline TikZ or a vector figure.}}
and the left Cauchy--Green (Finger) tensor is $\bt=\Ften\Ften^{\!\top}$. For an
incompressible motion $J\equiv1$; the discrete $J$ is monitored as a folding/quality
diagnostic.

\subsection{Neo-Hookean solid}
The equilibrium (non-relaxing) branch of the solid is an incompressible neo-Hookean
material,
\begin{equation}
\sig_{\!s}^{\infty} = \mu_s\,\dev\!\left(\Ften\Ften^{\!\top}\right)
   = \mu_s\!\left(\bt - \tfrac12\tr(\bt)\,\It\right)\quad(\text{2-D}),
\label{eq:neohookean}
\end{equation}
with shear modulus $\mu_s$ and elastic wave speed $c_s=\sqrt{\mu_s/\rho}$. In the purely
hyperelastic limit ($\tau\to\infty$, \cref{sec:constitutive}) this is the entire solid
stress; the general viscoelastic solid adds a relaxing branch.

\subsection{Level set, extrapolation, reinitialization}
The solid occupies $\{\phi\le0\}$. The level set is reconstructed from the reference map
each step (the material boundary is the image of the initial boundary under $\xiv$), and
the reference map is extrapolated into a narrow band outside the solid so that the
gradient in \cref{eq:F} and the interfacial forces have valid stencils. Optional
level-set reinitialization (PDE or fast-marching) restores a signed-distance band. The
advection of \cref{eq:xi-advect} is discretized by a switchable scheme (semi-Lagrangian
RK4 backtrace by default; central/WENO or a conservative flux form for
differentiability, \cref{sec:time}).
